Large language models are increasingly used to generate stories, synthetic route descriptions, and planning traces for embodied agents. Yet fluent spatial prose may still encode movement, landmarks, and perspective differently from human-authored spatial language. We test whether \llm stories about built-space movement exhibit a measurable \spatialgrammar. We build a balanced corpus of 114 human texts and 114 real \gptmini generations: 80 \rtr indoor route cases and 34 spatial \writingprompts prose-control cases. We extract transparent regex/count features covering length and generic style, spatial grammar, and narrative embellishment, then evaluate feature differences and logistic-regression author classifiers under stratified 5-fold cross-validation. Human and \llm texts are highly separable. Across the full corpus, all features reach 97.4\% accuracy, all non-length features reach 96.5\%, and spatial-only features reach 87.7\% with permutation $p=0.0099$. The \writingprompts control also remains separable: spatial-only features reach 83.8\% accuracy and all non-length features reach 94.1\%. The dominant \gptmini profile is first-person, sensory, formulaic, landmark-reusing, and less imperative than human route instructions. These results support a measurable \llm spatial-narrative fingerprint under the tested prompting setup, while stopping short of claims about stable cognitive maps or model-internal spatial representations.
